\documentclass[letterpaper,oneside,12pt,spanish]{book}

\usepackage[utf8]{inputenc}
\usepackage[spanish]{babel}
\usepackage[T1]{fontenc}
\usepackage{tocbibind} % Bibliografía en el indice
\usepackage{titlesec} % Posibilidad de editar los formatos de chapter
\usepackage{amsmath,amssymb,mathrsfs} % Matemáticas varias
\usepackage{hyperref} %Esto te hace un pequeño esquemita al lado
\usepackage{listings}
\usepackage{anysize} 
\usepackage{algorithmic} % Para pseudocódigo de algoritmos
\usepackage{algorithm}
\usepackage{appendix} % Para tener apéndice en vez de capítulos
\usepackage{verbatim}
\usepackage{multirow}
\usepackage{array}
\usepackage{pdflscape}




% --- Arreglos varios para la inclusion de imagenes
\usepackage[pdftex]{graphicx}
\usepackage{subfigure}
\usepackage{graphicx}
\usepackage[usenames,dvipsnames]{color}
\DeclareGraphicsExtensions{.png,.jpg,.pdf,.mps,.gif,.bmp}

% --- Para las dimensiones de los márgenes etc
% \frenchspacing \addtolength{\hoffset}{-1.5cm}
% \addtolength{\textwidth}{3cm} \addtolength{\voffset}{-2.5cm}
% \addtolength{\textheight}{4cm}
\marginsize{3.8cm}{2.5cm}{2cm}{2cm} 

\setlength{\parskip}{0.7em}

% --- Para el encabezado
\usepackage{fancyhdr}
\fancyhead[R]{UTFSM}\fancyhead[L]{} \fancyfoot[C]{\thepage}
\pagestyle{fancy}

% --- Para las tablas
\renewcommand{\tablename}{Tabla}
\renewcommand{\listtablename}{Índice de tablas}

% --- Para los apendices
\renewcommand{\appendixname}{Anexos}
%\renewcommand{\appendixtocname}{Anexos}
%\renewcommand{\appendixpagename}{Anexos}
% -------------------------------------------------------- %

\title{•}
\author{Paulina Andrea Silva Ghio}

\lstset{ %
language=C,                % choose the language of the code
basicstyle=\footnotesize,       % the size of the fonts that are used for the code
numbers=left,                   % where to put the line-numbers
numberstyle=\footnotesize,      % the size of the fonts that are used for the line-numbers
stepnumber=0,                   % the step between two line-numbers. If it's 1 each line 
                                % will be numbered
numbersep=5pt,                  % how far the line-numbers are from the code
backgroundcolor=\color{white},  % choose the background color. You must add \usepackage{color}
showspaces=false,               % show spaces adding particular underscores
showstringspaces=false,         % underline spaces within strings
showtabs=false,                 % show tabs within strings adding particular underscores
frame=single,	                % adds a frame around the code
tabsize=2,	                % sets default tabsize to 2 spaces
captionpos=b,                   % sets the caption-position to bottom
breaklines=true,                % sets automatic line breaking
breakatwhitespace=false,        % sets if automatic breaks should only happen at whitespace
title=\lstname,                 % show the filename of files included with \lstinputlisting;
                                % also try caption instead of title
escapeinside={\%*}{*)},         % if you want to add a comment within your code
morekeywords={*,...}            % if you want to add more keywords to the set
}

\begin{document}
\frontmatter
\begin{titlepage}

\begin{center}

\textsc{\Large Universidad Técnica Federico Santa María}\\
\textsc{\large Departamento de Informática}\\
\textsc{\large Valparaíso, Chile}\\[1.5cm]

% Upper part of the page
\includegraphics[width=0.3\textwidth]{figures/utfsm.jpg}\\[1cm]    

% Title
{ \huge Una Arquitectura de Referencia para Web Browsers: Hacia una unificación de Conceptos de Seguridad }\\[2cm]

% Author and supervisor
\text{\Large Paulina Andrea Silva Ghio}\\[2cm]
\text{\large Memoria para optar al título de: Ingeniería Civil Informática}\\[3cm]
\text{\large Profesor Guía: Raúl Monge}\\
\text{\large Profesor Correferente: Javier Cañas}

\vfill

% Bottom of the page
{\large \today}

\end{center}

\end{titlepage}
 
\begin{flushright}  
\null \vspace{\stretch{1}}
\emph{Si sta come \\d'autunno \\sugli alberi \\le foglie \\ \textbf{Soldati} di Guiseppe Ungaretti (1918)} 
\end{flushright}
\newpage 
\chapter*{Agradecimientos}
\label{chap:agrad}

Con esta Memoria voy a dar por finalizado los largos 6 años de la carrera, que terminaron por ser casí 7 y algo más por el Magister. No me arrepiento de mis decisiones en la parte académica, pero a veces en lo personal creo que pudo haber sido de otra manera. Tuve muchos desafíos que realizar al estudiar para este trabajo, partiendo por mi nulo conocimiento en Seguridad y a veces la falta de motivación. Sin embargo, estoy acá en este momento gracias a esas decisiones y experiencias que he ganado en el intertanto. Gracias a todas las personas que me motivaron, en especial, a mi profesor guía Raúl Monge (UTFSM) y el profesor Eduardo Fernandez (Florida Atlantic University), quienes me mostraron un camino interesante a seguir.

Quiero agradecer a todos mis amigos de la Universidad, que son muchos, y aquellos que la abandonaron por un camino mejor o más interesante. Gracias por haberse presentado en mi vida y hacerla entretenida. Sin la risa y diversión en mi vida, probablemente sería una persona amargada. Gracias a todos por soportar mi malhumor y comprenderme cuando estoy angustiada. Gracias también por haberme ayudado academicamente, que si bien todos creen que soy matea nada tengo de eso; solo estudio demasiado y a veces no de la mejor manera.

Gracias Mamá y Papá por darme educación. Se que no es fácil pagar tanto dinero todos los meses, pero en verdad se los agradezco siempre. Gracias a que he tenido la oportunidad de estudiar en una Universidad, se cuán importante es la educación y que debo aprovecharla al máximo. Se que muchas veces tuvimos varios problemas mientras yo crecía, pero quiero hacerles saber que estoy agradecida de ser su hija. Gracias por cuidar y confiar en mi.

Gracias Panchomon por todo, eres una existencia muy importante en mi vida y espero que lo sepas. Gracias Micchan, por ser mi manta cuando estoy triste. Ustedes dos han sido los pilares más grandes que he tenido en estos años de universidad. Gracias Fernando, Samuel, José Miguel, Francisco, Matías y Andrés, por ser mis \textit{dudes} y soportarme siempre. Gracias Priscilla, Scarlet y Susana, por ser mis amigas más cercanas en la Universidad. Ko-chan, Yurie, Asumi, Geso y Motoki a ustedes también gracias por apoyarme siempre. Finalmente, ¡gracias a todos!, por tenerme una gran paciencia a lo largo de estos años, por levantarme el ánimo cuando no quiero dar más pasos hacia adelante, por inspirarme a ser mejor, y por estar en mi vida de alguna que otra forma. También muchas gracias al que lee este documento, en especial, por darse el tiempo de hacerlo. 
\newpage 
% Resumen a grandes rasgos de lo existente (resumen del estado del arte y resumen del trabajo realizado)

\section*{Resumen}
\label{chap:resumen}

El \textit{Web Browser} es una de las aplicaciones más usadas - \textit{killer app} - y también una de las primeras en aparecer en cuanto se creó el Internet (década de los 90). Por lo mismo, su nivel de madurez con respecto a otros desarrollos es significativo y permite asegurar ciertos niveles de confianza cuando otros usan un \textit{Web Browser} como cliente para sus Sistemas. 

Actualmente muchos desarrollos de software crean sistemas que están conectados a la Internet, pues permite agregar funcionalidades al sistema y facilidades para sus \textit{Stakeholders}. Esto lleva a depender de un cliente web, cómo un \textit{Web Browser}, que permite el acceso a los servicios, datos u operaciones que el sistema entrega. Sin embargo, la Internet influye en la superficie de ataque del nuevo sistema, y lamentablemente tanto Stakeholders como muchos desarrolladores no están al tanto de los riesgos a los que se exponen.

Al tener sistemas que se interconectan con el \textit{Web Browser}, Stakeholder como Desarrolladores deben estar al tanto de los posibles riesgos que podrían enfrentar. La falta de educación de seguridad en los desarrolladores de software de un proyecto, la poca y dispersa documentación de cada navegador (así como su estandarización), podría llegar a ser un flanco débil en el desarrollo de grandes Architecturas que dependen del \textit{Browser} para realizar sus servicios. Una Arquitectura de Referencia del \textit{Web Browser}, utilizando Patrones Arquitecturales, podría ser una base para el entendimiento de los mecanismos de seguridad y su Arquitectura, que interactua con un sistema Web mayor. Ésto mismo, entregaría una unificación de ideas y terminología, al dar una mirada holística sin tener en cuenta detalles de implementación tanto del \textit{Browser} como el sistema con el que interactua.

En esta memoria presentada al Departamento de Informática (DI) de la UTFSM\footnote{Universidad Técnica Federico Santa María} Casa Central, incursionará en el ámbito de la seguridad del \textit{Web Browser}, tiene como objetivo el obtener documentos semi-formales que servirán como herramientas a personas que desarrollen Software y hagan un fuerte uso del navegador para las actividades del sistema desarrollado.

\section*{Abstract}
\label{chap:abstract}

The \textit{Web Browser} is known as one of the most used applications - or \textit{killer app} - and also was the first introduced when the Internet was created (1990s). Which is why, its significant maturity level is above in comparison with other developements and can assure a certain level of \textit{trust} whenever it is used as a client with other systems.

Currently a lot of software developments create systems that are connected to the Internet, which allows to add functionality within a system and facilities to their \textit{Stakeholders}. This leads to depend in a \textit{web client}, as the \textit{Web Browser}, which allows access to services, data or operations that the system delivers. Nevertheless, the Internet influences the attack surface of the new system, and unfortunately many stakeholders and developers are not aware of the risks they are exposed.

Having systems which are interconnected with the \textit{Web Browser}, Stakeholder and Developers should be aware of the potential risks they could face. The lack of Security Education in Software developers of a project, the low and scattered documentation of each browser (and standardization), could become a great flaw in big architectural developments which depends on the browser to do their services. A Reference Architecture of the \textit{Web Browser}, using Architectural Patterns, could be a base for understanding the security mechanisms and its architecture, which interacts with a bigger web system. This would give an unification of ideas and terminology, giving a holistic view regardless the implementation details for both the browser and the system it communicates to.

This work presented to the Departamento de Informática (DI) of the UTFSM\footnote{Universidad Técnica Federico Santa María} Casa Central, will seek within the scope of \textit{Web Browser} Security, has the objective or goal to obtain semi-formal documents that can help as tools to software developers who make a strong use of a web browser within the activities of the systems they are building.




 
\markboth{}{}
\newpage
\tableofcontents 
\newpage
\listoffigures
\newpage
\listoftables
\newpage	
\renewcommand{\sectionmark}[1]{\markright{\thesection\ #1}}

\mainmatter

%Intro CHECKED!!


\chapter{Introduction}
\label{chap1:intro}


\newpage
\chapter{Theoric Framework - Web Technologies}
\label{chap2:MT}


\newpage
\chapter{Browser's Threats}
\label{chap3:MT}


\newpage
\chapter{State of the Art}
\label{chap4:EA}


\newpage
\chapter{Towards a definition of a Security Reference Architecture}
\label{chap5:ArqRefWB}


\newpage
\chapter{Misuse Patterns}
\label{chap6:Misuse}


\newpage
\chapter{Conclusions}
\label{chap7:concl}

\newpage
\chapter{Glossary}
\label{Glosario}


%\backmatter
%\addcontentsline{toc}{chapter}{References}
\bibliography{refTodas}
\bibliographystyle{IEEEtran}

\appendix
%\input{chapters/anexos.tex}
%\input{chapters/anexo2.tex}
\end{document}